\chapter{Conclusiones generales} \label{cap:conclusiones_generales}

El ensayo a escala reducida arrastra una limitación conocida desde los orígenes
de la disciplina: la imposibilidad de satisfacer simultáneamente la semejanza de
todos los números adimensionales que gobiernan el flujo. En la resistencia al
avance, un modelo ensayado en el mismo fluido no puede igualar a la vez el número
de Froude y el de Reynolds del buque: al reducir la escala manteniendo el mismo
fluido y la misma gravedad, conservar el número de Froude obliga a operar a un número
de Reynolds mucho menor que el del buque. En la propulsión, la
misma tensión reaparece entre $Re$ y el coeficiente de avance $J$. Esta semejanza
dinámica parcial es la que introduce los efectos de escala, y su corrección por
vías indirectas constituye el punto de partida de la disciplina, muy anterior a
esta tesis.

Sobre esa dificultad se construyó, mucho antes del CFD, un conjunto de métodos
aproximados. La estrategia fundacional de Froude, respetar la semejanza de Froude
y corregir analíticamente la componente viscosa, dio lugar al factor de forma de
Hughes, al método de extrapolación de Prohaska, a la línea de correlación ITTC-57
y al procedimiento ITTC-78 para hélices, junto con diversas correlaciones
empíricas de corrección de escala. Cada uno de estos métodos resuelve una parte
del problema y arrastra sus limitaciones, porque la corrección viscosa es una
aproximación de un efecto que no puede medirse directamente a escala de modelo. La
línea friccional ITTC no es, en ese sentido, la causa de las dificultades sino una
de las herramientas de ese marco aproximado; su insuficiencia, cuando aparece, es un
síntoma de la aproximación subyacente y no su raíz.

El CFD modifica este panorama, porque permite eludir la limitación de la
semejanza: una simulación puede resolverse directamente al Reynolds de escala
buque, o recurrir a esquemas de escalamiento viscoso, sin la restricción que ata
al ensayo físico. Esa misma virtud, sin embargo, trae aparejada una complicación
que comparte con el ensayo a escala de buque: la ausencia de datos
experimentales con los cuales contrastar el resultado a esa escala. De ahí que el
trabajo se apoye en una combinación deliberada de EFD y CFD, en la que el ensayo
valida la herramienta numérica en el régimen accesible y el CFD, una vez validado,
extiende el análisis hasta la escala buque y descompone las fuerzas en sus
componentes de fricción y de presión, algo fuera del alcance de la medición
integral. Sobre esa base, la tesis contribuye por dos vías complementarias de la
cadena propulsiva: la resistencia al avance, a través del factor de forma $(1+k)$,
y la propulsión, a través del comportamiento de la hélice en aguas abiertas.

\section{Síntesis del Eje 1: resistencia al avance en cascos de baja esbeltez}

El primer eje abordó un tipo de buque escasamente estudiado en la literatura de
hidrodinámica naval: los pesqueros de baja esbeltez ($L/B < 4$), cuyos efectos de
escala en resistencia al avance no habían sido caracterizados con la profundidad
con que sí lo fueron los buques mercantes convencionales. El estudio experimental
del pesquero \textit{P1} (escala 1:20, tres calados) y del casco de referencia
\textit{KCS} (escala 1:79) permitió construir una base de datos de resistencia con
incertidumbres compatibles con las recomendaciones de la ITTC, cuantificadas tanto
por combinación cuadrática (RSS) como mediante el ajuste ponderado de York sobre el
diagrama de Prohaska. Esa base reveló ya una primera asimetría: mientras el
\textit{KCS} satisface con fidelidad los supuestos del método de Prohaska, con sus
puntos alineados sobre la recta teórica y un valor $(1+k) = 1{.}052$ de
incertidumbre relativa $\pm 2{.}3\%$, el \textit{P1} exhibe una curvatura
sistemática a bajos números de Froude y una incertidumbre del ajuste mucho mayor
($8$--$11\%$), señal de que los métodos clásicos se cumplen con menor holgura en un
casco de baja esbeltez.

El aporte central de este eje surgió al extender el análisis a escala buque
mediante CFD en configuración de doble cuerpo (DB). A diferencia del \textit{KCS},
cuyo factor de forma se estabiliza al aumentar el número de Reynolds tal como
predice la teoría clásica, el $(1+k)$ del \textit{P1} crece de manera continua y no
asintótica hasta la escala 1:1. La diferencia entre bajo y alto $Re$ alcanza
aproximadamente un $40\%$ (frente a cerca del $8\%$ en el \textit{KCS}), y el valor
de escala buque $(1+k) \approx 1{.}72$ ($k_S \approx 0{.}72$) resulta casi un
$50\%$ superior al del \textit{KCS}. Un efecto de escala de esta magnitud sobre el
factor de forma es sustancialmente mayor que el observado en los cascos de
referencia caracterizados en el marco del comité de resistencia de la ITTC, y no
tiene precedente documentado para un casco desnudo. Este comportamiento no es una
singularidad del \textit{P1}: los tres pesqueros analizados (\textit{P1},
\textit{P2} y \textit{P3}, con $L/B$ entre $2{.}85$ y $3{.}53$) reproducen la misma
tendencia creciente, lo que la establece como un rasgo estructural de los cascos
con $L/B < 4$.

El mecanismo físico responsable se aisló mediante la descomposición de la
resistencia viscosa que solo el CFD DB hace posible. El crecimiento sostenido de
$(1+k)$ se debe a la persistencia de la componente de presión viscosa $C_{PV}$,
originada en un déficit de recuperación de presión en la popa llena del casco y en
estructuras vorticosas longitudinales desprendidas de la quilla y el pantoque, que
se mantienen prácticamente invariantes con la escala. Al disminuir $C_{F0}$ con el
Reynolds sin que $C_{PV}$ lo acompañe, el cociente que define el factor de forma no
puede estabilizarse. Los análisis complementarios permitieron descartar
explicaciones alternativas: la proa bulbo modifica $(1+k)$ pero no explica ni la
curvatura del diagrama ni la ausencia de asíntota, que persisten en la variante sin
bulbo; la popa espejo sumergida, tratada con el método de dos factores de forma
$2k$, aporta $k_{tr} = 0{.}025$, apenas un $5\%$ del efecto de escala total
($k_S - k_M = 0{.}522$); y la sustitución de la línea ITTC-57 por líneas de fricción
numéricas atenúa la dependencia con $Re$ pero no la elimina. El efecto dominante es,
por lo tanto, la evolución de $C_{PV}$ gobernada por la geometría de baja esbeltez.

Como consecuencia directa, las correlaciones empíricas navales de corrección de
escala fallan en este régimen: la fórmula de \citet{gomezgarcia}, por ejemplo,
subestima el efecto de escala del \textit{P1} en más de un orden de magnitud,
porque fue calibrada para cascos con $L/B$ entre $6{.}98$ y $9{.}07$, muy lejos del
dominio de los pesqueros. En cambio, las correlaciones aerodinámicas de factor de
forma para cuerpos de revolución de baja esbeltez (\citealp{raymer2019aircraft,
roskam1987airplane}) reproducen con notable fidelidad los valores de CFD DB de los
tres pesqueros, ubicándolos en una región coherente con su $L/B$. Este hallazgo
abre la posibilidad de desarrollar una formulación empírica específica de $(1+k)$
en función de $L/B$ para cascos de baja esbeltez, que hoy carecen de una
herramienta de estimación confiable.

\section{Síntesis del Eje 2: efectos de escala en la hélice}

El segundo eje trasladó la misma dificultad de fondo, la imposibilidad de conservar
simultáneamente $Re$ y $J$ en el ensayo, al comportamiento de la hélice en aguas
abiertas, caracterizando dos mecanismos físicamente independientes de efecto de
escala. Ambos comparten el punto ciego del procedimiento ITTC-78: la suposición de
que la corrección de modelo a buque es puramente friccional ($\Delta K_{Tp} = 0$).

El primero es el efecto de escala entre condiciones completamente turbulentas (de
modelo a buque, ambas en régimen turbulento pleno), estudiado sobre la hélice
Wageningen B4-60 mediante un barrido de diámetros de $D = 0{.}14$ a $8{.}0$~m. Al
aumentar el Reynolds, el empuje crece ($\Delta K_T = +4{.}7\%$) y el par disminuye
($-4{.}1\%$), con una ganancia neta de eficiencia en aguas abiertas de $+9{.}1\%$.
La descomposición directa de fuerzas mostró que ese aumento de empuje está dominado
por la presión: la fracción $\Delta K_{Tp}/\Delta K_T$ se estabiliza en torno al
$76$--$77\%$, de modo que apenas una cuarta parte proviene de la reducción de
fricción. El análisis del campo de presión localizó esta redistribución sobre el
lado de succión de la mitad externa de la pala, intensificándose hacia la punta. El
procedimiento ITTC-78, al corregir solo fricción, predice una ganancia de eficiencia
de apenas $+1{.}1\%$ frente al $+9{.}1\%$ del CFD, una diferencia de unos $8$ puntos
porcentuales que corresponde íntegramente a la componente de presión ignorada.

El segundo mecanismo es la transición laminar-turbulenta sobre las palas a escala de
modelo, estudiada sobre el propulsor \textit{P206}, ensayado en túnel de cavitación
a cinco números de Reynolds de cuerda. Los ensayos revelan una joroba transicional
en $K_T$ y $K_Q$ cuya amplitud, medida como la diferencia entre los coeficientes
experimentales y la predicción completamente turbulenta, crece al descargar el
propulsor, desde un $7{.}1\%$ a $J = 0{.}30$ hasta un $16{.}3\%$ a $J = 0{.}60$. El
modelo completamente turbulento no reproduce este comportamiento en ningún caso, y
el modelo de transición ($\gamma$-$\widetilde{Re}_{\theta t}$), verificado en dos
solvers independientes (OpenFOAM y StarCCM+), lo captura solo parcialmente: ningún
valor único de intensidad turbulenta reproduce la forma de la curva en todo el rango
de operación, lo que impide tratar ese parámetro como una propiedad transferible.

No fue posible reproducir con exactitud los valores experimentales de $K_T$ y $K_Q$ en el régimen
transicional con ningún valor de intensidad turbulenta ni con ninguno de los dos
solvers empleados. Sin embargo, la descomposición directa de las fuerzas mostró que
el desplazamiento de la curva completamente turbulenta hacia los valores
experimentales es de origen mayoritariamente de presión --del orden del $90\%$--: la
componente friccional aporta la fracción menor, y el par incluso aumenta levemente
con el modelo de transición, cuando una reducción puramente friccional lo disminuiría.
La limitación remanente no es, entonces, la de identificar el mecanismo, sino la de que
ningún nivel único de intensidad turbulenta reproduce la joroba experimental en todo el
rango de operación. Al propagar esta física a
través del ITTC-78 se constató un doble problema: un sesgo sistemático de
subestimación de hasta el $8{.}5\%$, y, más grave, que el resultado de la
extrapolación depende críticamente de en qué velocidad de giro se realizó el único
ensayo de túnel disponible, de modo que si esa condición cae dentro de la joroba el
error puede duplicarse o triplicarse respecto de un punto vecino a Reynolds apenas
distinto.

La fracción de presión del efecto de escala completamente turbulento del
\textit{P206} se midió de forma directa, como el cociente $\Delta K_{Tp}/\Delta K_T$
obtenido al descomponer las fuerzas sobre la pala en sus componentes de presión y
fricción entre la condición de modelo y la de escala buque, ambas completamente
turbulentas, resultando $f_p \approx 0{.}90$--$0{.}92$. El valor es mayor que el de la
B4-60 ($0{.}70$--$0{.}77$) y prácticamente coincide con el de la hélice geométricamente
más parecida al \textit{P206} (el P1727, $\approx 0{.}88$), lo que sugiere que la
elevada fracción de presión está gobernada por el área expandida reducida; establecer
esa dependencia geométrica con generalidad requiere el mapeo sistemático que se plantea
como trabajo futuro.

A diferencia del eje de resistencia, donde la evidencia acumulada sobre tres cascos
permitió proponer una generalización geométrica del comportamiento en función de
$L/B$, el eje propulsivo deja el problema mucho más abierto. Los resultados aquí
obtenidos se apoyan en solo dos hélices estudiadas en detalle, la B4-60 y el
\textit{P206}, y en una tercera tomada como referencia externa de la bibliografía
reciente. La dependencia de la fracción de presión y de la joroba transicional con
los parámetros geométricos de la pala (número de palas, razón de áreas $A_E/A_0$,
skew, distribución de cuerdas) apenas comienza a esbozarse, y su generalización,
análoga a la lograda para el factor de forma con $L/B$, requiere una base de
geometrías considerablemente más amplia. El campo de los efectos de escala en
hélices permanece, en consecuencia, sensiblemente más abierto que el de la
resistencia.

\section{El hilo conductor: la componente de presión del efecto de escala}

Leídos en conjunto, los dos ejes convergen en una misma conclusión física. Frente a
la imposibilidad de la semejanza dinámica completa, los procedimientos vigentes
corrigen la componente viscosa suponiendo que el efecto de escala es esencialmente
friccional. Esta tesis muestra, en dos contextos independientes de la cadena
propulsiva, que ese supuesto es incompleto: en la resistencia, la incapacidad de
$(1+k)$ para estabilizarse con la escala se explica por una componente de presión
viscosa que persiste con el Reynolds; en la propulsión, el efecto de escala entre
condiciones completamente turbulentas está dominado por la redistribución de presión
sobre la pala, y la joroba transicional apunta, con la cautela ya señalada, en la
misma dirección. En ambos casos, el mecanismo que gobierna el escalado no es el que
los métodos estándar corrigen: la parte friccional captura la fracción menor del
fenómeno y omite la mayor. Haber cuantificado esa componente de presión constituye la
contribución transversal del trabajo.

\section{Impacto sobre la predicción de potencia del buque}

La consecuencia práctica de estos resultados se materializa al estimar la potencia
del \textit{P1} a escala buque. En el eje de resistencia, la comparación entre el
método ITTC estándar (que asume $k_S = k_M$) y un método de doble factor de forma
basado en CFD DB (con $k_M$ y $k_S$ determinados numéricamente en cada escala)
arroja diferencias del $13$ al $17\%$ en la potencia efectiva, gobernadas
directamente por el salto $k_S \gg k_M$ que constituye el resultado central del
primer eje. A modo de referencia, ignorar por completo el factor de forma ($k = 0$)
produce una sobrestimación de solo un $5\%$, lo que evidencia que la fuente de
discrepancia no es el procedimiento en sí, sino la magnitud del efecto de escala del
factor de forma, mucho mayor que el que asumen implícitamente los métodos vigentes.

En el eje propulsivo, la subestimación de la eficiencia en aguas abiertas, de unos
$8$ puntos porcentuales por el efecto de escala completamente turbulento (la
diferencia entre el $+9{.}1\%$ del CFD y el $+1{.}1\%$ del ITTC-78) y de hasta un
$8{.}5\%$ por el sesgo transicional del procedimiento, actúa sobre el otro extremo
del cálculo de potencia, degradando la predicción del rendimiento propulsivo. Dado
que resistencia y propulsión se combinan en la estimación de la potencia instalada,
sus contribuciones no se cancelan sino que se acumulan. Discrepancias del orden del
$15\%$ en potencia son suficientemente significativas desde el punto de vista del
diseño como para justificar su validación mediante pruebas de mar y para
desaconsejar la aplicación acrítica de los procedimientos estándar.

\section{El aporte metodológico: la sinergia EFD-CFD}

El resultado transversal de la tesis, tanto en resistencia como en propulsión, solo
fue posible gracias a la combinación deliberada de experimentación y simulación, en
la que cada herramienta suple las limitaciones de la otra. El CFD permite eludir la
restricción de la semejanza y alcanzar la escala buque, pero, al igual que el ensayo
a escala de buque, carece de un dato experimental con el cual contrastarse a
esa escala; el ensayo físico, a su vez, solo accede al régimen de modelo. La
validación cruzada entre ambos, en el régimen donde el ensayo es posible, no garantiza
por sí sola la exactitud del resultado a escala buque --que permanece sin un dato
experimental que lo confirme--, pero sí respalda las tendencias que la simulación
predice al cambiar de escala: la evolución de las distribuciones de presión, de las
fuerzas y de los regímenes de flujo. Es esto lo que permite comprender qué mecanismos
pueden estar gobernando el comportamiento a escala real.

La experimentación aportó la base validada e imprescindible: curvas de resistencia y
de aguas abiertas con incertidumbre cuantificada, que anclan el análisis en datos
reales. El acuerdo entre los ensayos del LabHiNO y las referencias internacionales,
las mediciones del \textit{KCS} frente a los datos de distintos laboratorios y las de
la B4-60 frente a MARIN, confirmó la calidad de la campaña experimental y habilitó el
uso del canal como fuente confiable de validación.

El CFD, a su vez, aportó lo que el ensayo no puede: por un lado, el acceso a escalas
de Reynolds inalcanzables en un canal de dimensiones finitas, que permitió observar
la evolución de $(1+k)$ y de los coeficientes de la hélice hasta la escala buque; por
otro, la descomposición de las fuerzas en sus componentes de fricción y de presión,
sin la cual el mecanismo físico del efecto de escala habría permanecido oculto tras
una fuerza integral. Fue esta descomposición la que reveló el papel dominante de
$C_{PV}$ en el factor de forma y la que permitió cuantificar la fracción de presión
$f_p$ de la hélice.

La sinergia es, además, específica de cada problema. En el factor de forma, la
combinación EFD-CFD hizo posible el método de doble cuerpo, que determina $k_M$ y
$k_S$ de forma independiente y resulta más robusto que el método de Prohaska
precisamente en los cascos donde este último se vuelve menos confiable. En
propulsión, la combinación permitió discriminar el efecto de escala completamente
turbulento del mecanismo transicional, que un modelo totalmente turbulento no puede
capturar y que solo el CFD de transición, anclado en datos de túnel, reproduce.

\section{Aportes principales}

Los aportes originales de esta tesis pueden resumirse en los siguientes puntos:

\begin{itemize}
    \item Se caracterizó en profundidad, por primera vez sobre datos propios y para
    tres cascos distintos, el efecto de escala del factor de forma en pesqueros de
    baja esbeltez ($L/B < 4$), un tipo de buque escasamente estudiado, mostrando que
    $(1+k)$ crece de manera continua y no asintótica con el número de Reynolds hasta
    la escala buque, con una magnitud muy superior a la de los cascos de referencia
    caracterizados en el marco del comité de resistencia de la ITTC.

    \item Se identificó a la componente de presión viscosa $C_{PV}$ de la popa como la
    causa física dominante de ese comportamiento, y se cuantificó y descartó, mediante
    los análisis de proa bulbo, popa espejo sumergida ($2k$) y línea de fricción, el
    peso de las causas alternativas.

    \item Se mostró la inaplicabilidad de las correlaciones navales de factor de forma
    a esta clase de cascos y se propuso el uso de correlaciones aerodinámicas de
    cuerpos de revolución en función de $L/B$ como marco más apropiado, abriendo la vía
    a una formulación empírica específica.

    \item Se cuantificó, para la hélice B4-60, que la componente de presión domina el
    efecto de escala completamente turbulento sobre el empuje ($\sim 76$--$77\%$).

    \item Se caracterizó, experimental y numéricamente, el comportamiento del propulsor
    \textit{P206} en régimen transicional, describiendo la joroba en $K_T$ y $K_Q$ y
    poniendo en evidencia el problema que supone extrapolar las prestaciones desde ese
    régimen, donde el resultado depende críticamente del número de Reynolds al que se
    realizó el ensayo.

    \item Se analizó, por separado, el efecto de escala completamente turbulento del
    \textit{P206}, cuantificando su fracción de presión en $f_p \approx 0{.}90$--$0{.}92$
    mediante la descomposición directa de fuerzas, un valor muy superior al de las
    hélices convencionales de área moderada.

    \item Se evaluó el impacto de estos efectos sobre la predicción de potencia,
    mostrando diferencias del orden del $13$--$17\%$ en potencia efectiva por el factor
    de forma y de hasta el $8{.}5\%$ en eficiencia propulsiva por el sesgo del ITTC-78,
    y se propusieron métodos de doble $k$ y correcciones de fricción y presión que los
    incorporan.

    \item Se consolidó una metodología combinada EFD-CFD, validada frente a
    referencias internacionales, capaz de abordar problemas de escala que ni la
    experimentación ni la simulación resuelven por separado.
\end{itemize}

\section{Líneas futuras de investigación}

El trabajo realizado deja abiertas varias direcciones de continuación, con un grado
de madurez distinto en cada eje. En resistencia, la caracterización alcanzada permite
plantear objetivos de consolidación; en propulsión, el campo se encuentra
considerablemente más abierto.

\begin{itemize}
    \item El desarrollo y la calibración de una correlación empírica de $(1+k)$ en
    función de $L/B$ para cascos pesqueros, específicamente en el rango $L/B < 4$, a
    partir de una base ampliada de geometrías, que permita estimar $k_S$ sin recurrir
    a simulaciones CFD a escala buque.

    \item La validación de las predicciones de potencia mediante pruebas de mar sobre
    buques en escala 1:1, dada la magnitud de las discrepancias frente a los métodos estándar.

    \item En el terreno propulsivo, la búsqueda de una generalización del efecto de
    escala y de la fracción de presión en función de la geometría de la pala (número de
    palas, razón de áreas $A_E/A_0$, skew, distribución de cuerdas), a partir de un
    conjunto de hélices mucho más amplio que las estudiadas aquí. Es la contrapartida
    propulsiva de la generalización lograda con $L/B$ para la resistencia, y permanece
    esencialmente pendiente.

    \item La exploración de la introducción deliberada de turbulencia sobre las palas,
    mediante turbuladores o rugosidad controlada, como estrategia para forzar un régimen
    completamente turbulento a escala de modelo y evitar así la joroba transicional
    observada tanto en los ensayos (EFD) como en el modelo de transición (SST-LM). De
    confirmarse, esta vía volvería más predecible el ensayo de aguas abiertas y su extrapolación,
    al ubicar el punto de operación fuera del rango ambiguo.

    \item La determinación de la forma funcional $f_p(Re)$ de la fracción de presión, y
    la incorporación de las correcciones de presión, en resistencia y en propulsión, a
    los procedimientos de predicción de potencia, hacia un esquema de extrapolación que
    no suponga que el efecto de escala es puramente friccional.

    \item La extensión del análisis a la condición autopropulsada completa, que integre
    la interacción casco-hélice y permita evaluar el impacto conjunto de ambos ejes
    sobre el punto de operación real del buque.
\end{itemize}

En conjunto, esta tesis aporta en un terreno poco explorado. En resistencia,
caracteriza el efecto de escala del factor de forma en pesqueros de baja esbeltez y
lo explica por una componente de presión que los métodos estándar no contemplan,
dejando al alcance una generalización geométrica en función de $L/B$. En propulsión,
cuantifica el peso de la presión en el efecto de escala de la hélice y describe la
joroba transicional, abriendo un terreno donde el mapeo del
comportamiento en función de la geometría de la pala queda por delante. En ambos
frentes, la combinación de experimentación y simulación numérica se muestra como la
herramienta adecuada: el ensayo aporta la validación en el régimen accesible, el CFD
el acceso a la escala real y a la física subyacente, y solo su uso conjunto permite
cuantificar la componente de presión del efecto de escala que los procedimientos
vigentes, apoyados en una hipótesis puramente friccional, dejan fuera. El eje de
resistencia queda así caracterizado con una generalización geométrica al alcance;
el de propulsión, con un campo todavía ancho por delante.